\section{未成年人犯罪案件诉讼程序}
\begin{enumerate}
  \item \zd{保护和教育}。
  \item 未成年人扩大解释。未成年被害人、证人。
  \item \zd{依附关系}，特别程序没有特别规定的，适用普通程序。
  \item 诉讼时 = 侦查阶段 + 审查起诉 + 审判阶段。
  \item 方针与原则。
        \begin{enumerate}
          \item 教育为主惩罚为辅。法庭辩论结束后，可以法制教育；宣判之后应当进行\zd{法治教育}。
          \item 最有利于未成年人与保护未成年人权利原则。防止二次伤害。
          \item 尽可能\zd{分案处理}。例外：如未成年人是首要分子。执行阶段一定是要分案关押。
          \item 不公开审理原则和保密原则。\zd{审理中}。
          \item 全面调查。
          \item 社会参与原则。
        \end{enumerate}
  \item 办案主体专门化，少年法庭。
        \begin{enumerate}
          \item 犯罪时不满 18 岁，立案时不满 20 周岁。应当。
          \item 立案时还没满 22 岁的在校大学生。可以。
        \end{enumerate}
  \item 立案审查年龄。12、14、16、18。\zd{没有查清无法确定}，做有利于被告人的推定，只能\zd{推定未满}。
  \item \zd{慎用}强制措施。应当主动听取辩护人意见。
  \item 义务教育没有完成，应当与教育行政部门互相配合保证他接受义务教育。
  \item 取保候审，应当指定合适的成年人作为他的保证人。
  \item \zd{附条件不起诉}。几乎每年都考。
        \begin{enumerate}
          \item 前提条件。小孩犯罪 456，1 年以下够起诉，悔过可以附条件。
          \item 被害人不服可以向上一级检察院申诉，不能自诉（不是彻底的不起诉）。
          \item 暂缓的不起诉，\zd{检察院}考验期 \zd{一年半载}，不计入审查起诉期限，先暂停。
          \item \zd{应当义务：市县、报告、守法、教育}。
          \item 处理结果：\yi 新罪、漏罪、违法、违规 \yjt 起诉；经受住考验 \yjt 应当不起诉。
        \end{enumerate}
  \item \zd{简易程序要三方同意}，绝不能用速裁程序。
  \item \zd{认罪认罚有异议，可以不签具结书，但不影响适用}。
  \item \zd{法代可以补充陈述}。
  \item \zd{公开宣判，但不组织旁听}。
  \item \zd{犯罪记录封存}。犯罪时。
        \begin{enumerate}
          \item 全过程及其周边材料；相对封存 \yjt 无罪的；一条龙封存；
          \item 不满 18 + 5 年以下以及免予刑事处罚。
          \item 再犯罪接受调查时，应当主动如实供述犯罪记录。
          \item \zd{申请出具无犯罪记录证明的，受理单位应当在 3 个工作日内出具}。
          \item \zd{解除封存：新罪，漏罪，超 5 年；成年后故意犯罪}。
        \end{enumerate}
\end{enumerate}
